\BourbakiExercises{习题}

\begin{enumerate}
\def\labelenumi{\arabic{enumi})}
\tightlist
\item
  设 \(\mathrm{K}\) 是一个交换域，\(\mathrm{A}\) 是一个 K-代数。若对每个扩张 \(\mathrm{L}\)（\(\mathrm{K}\) 的），\(\mathrm{A}_{(\mathrm{L})}\) -模 \(\mathrm{M}_{(\mathrm{L})}\) 都无根，则称 A-模 M 是绝对无根的。我们说代数 \(\mathrm{A}\) 是绝对无根的，如果 A-模 \(\mathrm{A}_s\) 是绝对无根的。
\end{enumerate}

\begin{enumerate}
\def\labelenumi{\alph{enumi})}
\item
  一个绝对无根的交换代数是可分的（V, §15, No.~2, 第 119 页, 定义 1）。
\item
  绝对无根的模的每个子模都是绝对无根的；绝对无根的模的一个直和是绝对无根的。
\item
  设 M 是一个有限生成的 A-模；则 M 绝对无根当且仅当存在一个完满域 P（K 的代数扩张），使得 \(\mathfrak{R}_{\mathrm{A}_{(\mathrm{P})}}(\mathrm{M}_{(\mathrm{P})}) = 0\)。
\item
  半单 A-模 M 绝对无根当且仅当属于 M 的支集的每个单模都绝对无根。
\item
  设 S 是一个单 A-模，D 是它的交换子。证明下述性质等价：
\end{enumerate}

\begin{enumerate}
\def\labelenumi{(\roman{enumi})}
\item
  A-模 S 是绝对无根的。
\item
  K-代数 \(\mathrm{D}\) 是绝对无根的。
\item
  \(\mathrm{D}\) 的中心是 \(\mathrm{K}\) 的一个可分扩张。
\end{enumerate}

\(f\) ) 一个半单代数绝对无根当且仅当它的每个单分量的中心都是 K 的一个可分扩张。

\(g\) ) 证明一个绝对无根且在 \(\mathbf{K}\) 上有限维的 A-模是绝对半单的。

\begin{enumerate}
\def\labelenumi{\arabic{enumi})}
\setcounter{enumi}{1}
\tightlist
\item
  设 \(A_{1}\) 与 \(A_{2}\) 是 K-代数，A 是代数 \(A_{1} \otimes_{K} A_{2}\)，\(M_{i}\) 是一个 \(A_{i}\) -模（对 i = 1, 2），M 是 A-模 \(M_{1} \otimes_{K} M_{2}\)。
\end{enumerate}

\begin{enumerate}
\def\labelenumi{\alph{enumi})}
\item
  假设 \(A_{1}\) -模 \(M_{1}\) 绝对无根。证明包含关系 \(\mathfrak{R}_{\mathrm{A}}(\mathrm{M}) \subset \mathrm{M}_{1} \otimes_{\mathrm{K}} \mathfrak{R}_{\mathrm{A}_{2}}(\mathrm{M}_{2})\)。
\item
  若模 \(M_{1}\) 与 \(M_{2}\) 都绝对无根，则 A-模 M 亦然。
\end{enumerate}

\begin{enumerate}
\def\labelenumi{\arabic{enumi})}
\setcounter{enumi}{2}
\item
  设 K 是一个交换域，A 与 B 是单 K-代数。假设 A 在其中心 Z 上秩有限，Z 是 K 的一个代数扩张，且 K-代数 B 绝对无根（习题 1）。证明环 \(A \otimes_{K} B\) 是绝对平坦的（VIII, 第 171 页, 习题 27；注意 \(A \otimes_{K} B\) 的每个元素都属于形如 \(A_{1} \otimes_{K} B\) 的子代数，其中 \(A_{1}\) 是 K 上秩有限的单代数）。
\item
  设 \(\mathrm{K}\) 是一个特征为 \(p > 0\) 的交换域，\(a\) 是 \(\mathrm{K} - \mathrm{K}^p\) 的一个元素。用 A 表示（有限维）K-代数 \(\mathrm{K}[\mathrm{X}] / ((\mathrm{X}^p - a)^2)\)。证明 A 的根 \(\mathfrak{r}\) 的平方为零，K-代数 \(\mathrm{A} / \mathfrak{r}\) 同构于 K 的根式扩张 \(\mathrm{L} = \mathrm{K}[\mathrm{X}] / (\mathrm{X}^{p} - a)\)，且 A 不包含任何同构于 L 的 K-子代数。
\item
  设 \(\mathrm{K}\) 是一个交换域，\(\mathrm{L}\) 是一个有限集。用 B 表示 VIII, 第 169 页习题 17 中构造的代数，把中心元素 \(1 - \sum_{\lambda} c_{\lambda} + \sum_{\lambda \neq \mu} r_{\lambda \mu}\)（\(\mathrm{B}\) 的）记作 \(\varepsilon\)，并用 A 表示代数 \(\mathrm{B} / \mathrm{B}\varepsilon\)。K-代数 A 是有限维的，它的根 \(\mathfrak{r}\) 的平方为零，且 \(\mathrm{A} / \mathfrak{r}\) 同构于 \(\mathrm{K}^{\mathrm{L}}\)。证明 A 不同构于两个非零代数的乘积。
\end{enumerate}

¶ 6) 设 K 是一个交换域，E 是 K 的一个扩张。证明下述性质等价：

\begin{enumerate}
\def\labelenumi{(\roman{enumi})}
\item
  对 K 的每个可分扩张 L，环 \(E_{(L)}\) 是半单的。
\item
  域 \(\mathrm{E}\) 是 \(\mathrm{K}\) 的一个有限次数可分扩张的纯不可分扩张。
\end{enumerate}

（假设 (i)，利用 V, §17, No.~2, 第 140 页的推论与 V, §2, 第 146 页的习题 4 证明 E 是 K 的一个代数扩张。然后证明 K 在 E 中的可分闭包 \(E_{s}\) 次数有限。为此，注意在相反的情形，环 \(\mathrm{E}_{(\mathrm{L})}\)（其中 L 是 K 的一个可分闭包）将包含 n 任意大的正交幂等元族 \((e_{1},\ldots,e_{n})\)。假设 (ii)，证明若 \(\mathrm{E}_{s(\mathrm{L})}\) 同构于 \(\prod_{j=1}^{r}C_{j}\)，则 \(\mathrm{E}_{(\mathrm{L})}\) 同构于 \(\prod_{j=1}^{r}D_{j}\)，其中 \(D_{j}\) 是 \(C_{j}\) 的一个纯不可分扩张。）

¶ 7) 设 E 是特征 p \textgreater{} 0 的交换域 K 的一个 p-根式扩张。假设环 \(E \otimes_{K} E\) 是诺特的。

\begin{enumerate}
\def\labelenumi{\alph{enumi})}
\item
  证明存在一个整数 \(k\)，使得 \(\mathrm{E} \subset \mathrm{K}^{p^{-k}}\)（用反证法推理：构造一个序列 \((x_{n})\)（在 \(\mathrm{E}\) 中），使得 \((x_{n} \otimes 1 - 1 \otimes x_{n})^{p^{n-1}} \neq 0\) 且 \((x_{n} \otimes 1 - 1 \otimes x_{n})^{p^n} = 0\)）。
\item
  设 \(\mathrm{F} = \mathrm{K}(\mathrm{E}^{p})\)，并设 \((x_{i})_{i \in \mathrm{I}}\) 是 E 在 F 上的一个 p-基（V, §13, No.~1, 第 98 页）。令 \(y_{i} = x_{i} \otimes 1\) 与 \(z_{i} = x_{i} \otimes 1 - 1 \otimes x_{i}\)。设 \(\Lambda\) 是 \(\mathbf{N}^{(I)}\) 的由族 \((\alpha_{i})_{i \in \mathrm{I}}\)（满足 \(\alpha_{i} < p\)，对每个 i）组成的子集。证明元素 \(\prod_{i \in I} y_{i}^{\alpha_{i}} z_{i}^{\beta_{i}}\)（对 \((\alpha_{i})\) , \((\beta_{i}) \in \Lambda\)）构成 F-向量空间 \(E \otimes_{F} E\) 的一个基。由此得出 I 是有限的（注意环 \(E \otimes_{F} E\) 是诺特的）。
\item
  由此得出 E 在 K 上次数有限（利用 V, §13, 第 170 页, 习题 1, b)）。
\end{enumerate}

¶ 8) 设 E 是交换域 K 的一个代数扩张。证明环 \(E \otimes_{K} E\) 是诺特的当且仅当 E 在 K 上次数有限（利用习题 6 的方法化归为纯不可分扩张的情形，然后应用习题 7）。

¶ 9) 设 K 是一个交换域，A 是一个 K-代数。假设对每个阿廷 K-代数 B，环 \(A \otimes_{K} B\) 都是阿廷的。证明 A 在 K 上有限维（通过注意 \(A/\Re(A)\) -模 \(\Re(A)^{k}/\Re(A)^{k+1}\) 有有限长度，化归为 A 半单的情形；然后利用习题 6 与 7）。

\begin{enumerate}
\def\labelenumi{\arabic{enumi})}
\setcounter{enumi}{9}
\tightlist
\item
  设 K 是一个交换环，A 是一个 K-代数，P 是一个 (A,A)-双模。A 的一个经 P 的扩张是一个三元组 \((B,i,p)\)，其中 B 是一个 K-代数，\(p:B\to A\) 是 K-代数的一个满同态，而 \(i:P\to B\) 是象为 Ker p 的单 K-线性同态，满足 \(i(p(b)x)=bi(x)\) 与 \(i(xp(b))=i(x)b\)（对 \(x\in P\) 与 \(b\in B\)）。从扩张 \((B,i,p)\) 到扩张 \((B',i',p')\) 的一个同构是 K-代数的一个同构 \(u:B\to B'\)，使得 \(p'\circ u=p\) 且 \(u\circ i=i'\)。
\end{enumerate}

\begin{enumerate}
\def\labelenumi{\alph{enumi})}
\item
  设 \((B, i, p)\) 是 A 的一个经 P 的扩张；i 的象是 B 的一个平方为零的双边理想。反之，若 \(p : B \to A\) 是 K-代数的一个满同态且其核 P 的平方为零，而 i 是 P 到 B 中的典范单射，则三元组 \((B, i, p)\) 是 A 的一个经 P 的扩张。
\item
  设 \(i_{0}: P \to A \oplus P\) 与 \(p_{0}: A \oplus P \to A\) 是典范映射。在 \(A \oplus P\) 上存在唯一的 K-代数结构，使得 \((A \oplus P, i_{0}, p_{0})\) 是 A 的一个经 P 的扩张。证明这个扩张的自同构群可以等同于从 A 到 P 的 K-导子所成的群。
\item
  设 \((B, i, p)\) 是 A 的一个经 P 的扩张。证明下述条件等价：
\end{enumerate}

\begin{enumerate}
\def\labelenumi{(\roman{enumi})}
\item
  扩张 \((\mathrm{B},i,p)\) 同构于扩张 \((\mathrm{A}\oplus \mathrm{P},i_0,p_0)\)。
\item
  存在 K-代数的一个同态 \(s: \mathrm{A} \to \mathrm{B}\)，使得 \(p \circ s = \operatorname{Id}_{\mathrm{A}}\)。
\item
  存在子代数 \(S\)（\(B\) 的），使得 \(B = S \oplus i(P)\)。
\end{enumerate}

若这些条件满足，则我们说扩张 \((B,i,p)\) 是平凡的。d) 设 \((B,i,p)\) 与 \((B',i',p')\) 是 A 的两个经 P 的扩张。设 C 是 \(B \times B'\) 的由元素 \((b,b')\)（满足 \(p(b) = p'(b')\)）组成的子代数，设 c 是 C 的由元素 \((i(x), -i'(x))\)（对 \(x \in P\)）组成的双边理想，并设 \(B''\) 是代数 C/c 且 \(\pi : C \to B''\) 是典范同态。设 \(i'' : P \to B''\) 与 \(p'' : B'' \to A\) 是由 \(i''(x) = \pi(i(x),0)\) 与 \(p''(\pi(b,b')) = p(b)\)（对 \(x \in P\) , \(b \in B\) , \(b' \in B'\)）定义的同态。证明 \((B'', i'', p'')\) 是 A 的一个经 P 的扩张；它称为 \((B,i,p)\) 与 \((B',i',p')\) 的融合和。证明我们这样就在 A 的经 P 的扩张的同构类所成的集合 \(\mathrm{Ex}_{\mathrm{K}}(\mathrm{A},\mathrm{P})\) 上定义了一个交换群律，其单位元是平凡扩张的类。此外，K 在 \(\mathrm{Ex}_{\mathrm{K}}(\mathrm{A},\mathrm{P})\) 上由 \(\lambda(\mathrm{B},i,p) = (\mathrm{B},\lambda^{-1}i,p)\)（对 \(\lambda \in K\) , \(\lambda \neq 0\)）给出的作用在 \(\mathrm{Ex}_{\mathrm{K}}(\mathrm{A},\mathrm{P})\) 上定义了一个 K-模结构。

\begin{enumerate}
\def\labelenumi{\alph{enumi})}
\setcounter{enumi}{4}
\tightlist
\item
  设 \(u: P \to P'\) 是 (A, A)-双模的一个同态，并设 \((B, i, p)\) 是 A 的一个经 P 的扩张。设 \(B'\) 是融合和 \(B \oplus_{P} P'\)（II, §1, 第 381 页, 习题 5）；对 \(b \in B\) , \(x' \in P'\)，把 \((b, x')\) 在 \(B'\) 中的类记作 \([b, x']\)。设 \(i': P' \to B'\) , \(p': B' \to A\) 与 \(v: B \to B'\) 是由下式定义的映射
\end{enumerate}

\[
i ^ {\prime} (x ^ {\prime}) = [ 0, x ^ {\prime} ], \qquad p ^ {\prime} ([ b, x ^ {\prime} ]) = p (b), \qquad v (b) = [ b, 0 ].
\]

证明在 \(B'\) 上存在唯一的代数结构，使得 \((\mathrm{B}', i', p')\) 是 A 的一个经 \(P'\) 的扩张且 v 是代数同态。我们说 \((\mathrm{B}', i', p')\) 是 A 的经 \(P'\) 的、由 \((\mathrm{B}, i, p)\) 经 u 得出的扩张；每个扩张 \((\mathrm{C}, j, q)\)（A 的经 \(P'\) 的扩张），若存在代数同态 \(w: B \to C\) 满足 \(q \circ w = p\) 与 \(w \circ i = j \circ v\)，则它同构于 \((\mathrm{B}', i', p')\)。

\(f\) ) 映射 \(\mathrm{Hom}_{(\mathrm{A},\mathrm{A})}(\mathrm{P},\mathrm{P}')\times \mathrm{Ex}_\mathrm{K}(\mathrm{A},\mathrm{P})\to \mathrm{Ex}_\mathrm{K}(\mathrm{A},\mathrm{P}')\) 把二元组 \((u,(B,i,p))\) 映到由 (B, \(i,p)\) 经 \(u\) 得出的扩张的类，它是 K-双线性的。\(g\) ) 设 \(f:\mathrm{A}'\to \mathrm{A}\) 是一个 K-代数同态，并设 (B, \(i,p)\) 是 A 的一个经 P 的扩张。类似地定义纤维积 \(\mathrm{B}'' = \mathrm{B}\times_{\mathrm{A}}\mathrm{A}'\) 上的代数结构以及扩张 \((\mathrm{B}'' , i'', p'')\)（\(\mathrm{A}'\) 的经 P 的扩张），它称为由 (B, \(i,p)\) 经 \(f\) 得出的扩张。

\begin{enumerate}
\def\labelenumi{\alph{enumi})}
\setcounter{enumi}{7}
\tightlist
\item
  再假设代数 A 是交换的。证明 A 的经 P 的、使得代数 B 为交换的扩张 \((B,i,p)\) 的类构成 \(\operatorname{Ex}_{\mathrm{K}}(\mathrm{A},\mathrm{P})\) 的一个 K-子模，我们把它记作 \(\operatorname{Ex}_{\mathrm{K}}^{c}(\mathrm{A},\mathrm{P})\)。
\end{enumerate}

\begin{enumerate}
\def\labelenumi{\arabic{enumi})}
\setcounter{enumi}{10}
\tightlist
\item
  \begin{enumerate}
  \def\labelenumii{\alph{enumii})}
  \tightlist
  \item
    设 f 是 \(\mathrm{C}^{2}(\mathrm{A},\mathrm{P})\) 的一个元素（VIII, 第 239 页）。证明我们在 K-模 \(A \oplus P\) 上通过令 \((a,x)(a',x') = (aa',ax' + xa' + f(a,a'))\) 来定义一个代数结构。典范单射 \(i: P \to A \oplus P\) 与典范满射 \(p: A \oplus P \to A\) 是环同态，且三元组 \((\mathrm{A} \oplus \mathrm{P}, i, p)\) 是 A 的一个经 P 的扩张（注意我们有 \(f(a,1) = af(1,1)\) 与 \(f(1,a) = f(1,1)a\)（对每个 \(a \in A\)））。
  \end{enumerate}
\end{enumerate}

\begin{enumerate}
\def\labelenumi{\alph{enumi})}
\setcounter{enumi}{1}
\item
  证明这个扩张的同构类只依赖于 f 在 \(\mathrm{H}^{2}(\mathrm{A},\mathrm{P})\) 中的类。于是 a) 中的构造定义了一个从 \(\mathrm{H}^{2}(\mathrm{A},\mathrm{P})\) 到 \(\mathrm{Ex}_{\mathrm{K}}(\mathrm{A},\mathrm{P})\) 的由作为 K-模的扩张为平凡的 A 的经 P 的扩张组成的 K-子模上的 K-模同构。
\item
  假设环 A 是交换的；用 \(\mathrm{C}^{2}(\mathrm{A},\mathrm{P})^{s}\) 表示 \(\mathrm{C}^{2}(\mathrm{A},\mathrm{P})\) 的由满足 \(f(x,y)=f(y,x)\)（对 A 中所有 x, y）的元素 f 组成的 K-子模，并用 \(\mathrm{H}^{2}(\mathrm{A},\mathrm{P})^{s}\) 表示 \(\mathrm{C}^{2}(\mathrm{A},\mathrm{P})^{s}\) 在 \(\mathrm{H}^{2}(\mathrm{A},\mathrm{P})\) 中的象。证明 b) 中的同构诱导一个从 \(\mathrm{H}^{2}(\mathrm{A},\mathrm{P})^{s}\) 到 \(\mathrm{Ex}_{\mathrm{K}}^{c}(\mathrm{A},\mathrm{P})\) 的由作为 K-模的扩张为平凡的交换的 A 的经 P 的扩张组成的 K-子模的同构。
\end{enumerate}

\begin{enumerate}
\def\labelenumi{\arabic{enumi})}
\setcounter{enumi}{11}
\tightlist
\item
  设 A 是一个交换环，J 是 A 的一个理想，B 是环 A/J，P 是一个 B-模。a) 证明正合列 \(0 \rightarrow J/J^{2} \rightarrow A/J^{2} \rightarrow B \rightarrow 0\) 定义了 A-代数 B 的一个经 \(J/J^{2}\) 的扩张（习题 10）。
\end{enumerate}

\begin{enumerate}
\def\labelenumi{\alph{enumi})}
\setcounter{enumi}{1}
\tightlist
\item
  设 P 是一个 B-模。证明把 u 映到由 a) 中的扩张经 u 得出的扩张的同态 \(\operatorname{Hom}_{\mathrm{B}}(\mathrm{J}/\mathrm{J}^{2},\mathrm{P}) \to \operatorname{Ex}_{\mathrm{A}}^{c}(\mathrm{B},\mathrm{P})\) 是一个同构。
\end{enumerate}

